\section{5차시 --- LED 3개로 신호등 만들기}
\label{sec:ch05}

본 차시는 LED 3개로 "화성 교통 신호등"을 만든다. 점등 순서를 설계
하는 \emph{시퀀스 사고}와, 단순한 무한 반복을 넘어선 \emph{모드 분기}
(주간 신호 / 야간 점멸 신호 / 화성 특별 신호)를 다룬다. 학습자는 사용자
정의 시나리오를 코드로 표현하는 첫 경험을 한다.

\subsection{미션 1 --- 신호등 LED 연결 점검}
\label{subsec:ch05-m1}
\missionmeta{LED 3개}{TRAFFIC}

\begin{storybox}
\sayares{알비! 신호등 판에 LED 3개를 꽂고 각각 켜지는지 먼저 확인해봐야 해!}
\sayalbi{삐릿- 맞아요! 빨강, 노랑, 초록을 각각 켜보며 정상 작동하는지 확인해봐요.}
\end{storybox}

\subsubsection*{학습 목표}
\begin{itemize}[leftmargin=1.4em]
  \item 빨강 LED(3번) → 2초 켜기 → 끄기
  \item 노랑 LED(2번) → 2초 켜기 → 끄기
  \item 초록 LED(1번) → 2초 켜기 → 끄기
  \item 순서대로 점등해 정상 작동 확인
\end{itemize}

\begin{labbox}{샘플 코드 --- 신호등 LED 점검}
\begin{minted}{python}
# 신호등 LED 점검
led_on(3, 1.0)   # 빨강
time.sleep(2)
led_off(3)
led_on(2, 1.0)   # 노랑
time.sleep(2)
led_off(2)
led_on(1, 1.0)   # 초록
time.sleep(2)
led_off(1)
\end{minted}
\end{labbox}

\subsection{미션 2 --- 화성 교통 신호등 작동!}
\label{subsec:ch05-m2}
\missionmeta{LED 3개}{TRAFFIC}

\begin{storybox}
\sayalbi{삐릿! 탐사 로봇들을 안전하게 시험 주행하려면 신호등 시스템이 필요해요!}
\sayares{초록→노랑→빨강 순서로 켜지는 화성 교통 신호등을 만들어보자!}
\end{storybox}

\subsubsection*{학습 목표}
\begin{itemize}[leftmargin=1.4em]
  \item 초록 LED 켜기 → 3초 → 끄기 (출발 신호)
  \item 노랑 LED 켜기 → 1초 → 끄기 (주의 신호)
  \item 빨강 LED 켜기 → 3초 → 끄기 (정지 신호)
  \item 3단계를 무한 반복하는 완전한 신호등
\end{itemize}

\begin{labbox}{샘플 코드 --- 화성 신호등}
\begin{minted}{python}
# 화성 신호등
while True:
    led_on(1, 1.0)   # 초록 출발
    time.sleep(3)
    led_off(1)
    led_on(2, 1.0)   # 노랑 주의
    time.sleep(1)
    led_off(2)
    led_on(3, 1.0)   # 빨강 정지
    time.sleep(3)
    led_off(3)
\end{minted}
\end{labbox}

\subsection{미션 3 --- 밤 신호등 점멸 신호}
\label{subsec:ch05-m3}
\missionmeta{LED 3개}{TRAFFIC}

\begin{storybox}
\sayares{알비, 화성에서 밤이 되면 신호등이 달라져야 하지 않을까?}
\sayalbi{삐릿- 좋은 관찰이에요! 밤 모드로 노랑 LED만 계속 깜빡이는 점멸 신호등을 만들어봐요!}
\end{storybox}

\subsubsection*{학습 목표}
\begin{itemize}[leftmargin=1.4em]
  \item 노랑 LED(2번)만 0.5초 간격으로 깜빡이기
  \item 빨강·초록은 꺼진 상태 유지
  \item 밤 모드 점멸 신호 무한 반복
\end{itemize}

\begin{labbox}{샘플 코드 --- 밤 신호등 점멸}
\begin{minted}{python}
# 밤 신호등 점멸
led_off(1); led_off(3)   # 빨강·초록 끄기
while True:
    led_on(2, 1.0)   # 노랑 켜기
    time.sleep(0.5)
    led_off(2)       # 노랑 끄기
    time.sleep(0.5)
\end{minted}
\end{labbox}

\subsection{미션 4 --- 나만의 화성 특별 신호}
\label{subsec:ch05-m4}
\missionmeta{LED 3개}{TRAFFIC}

\begin{storybox}
\sayares{알비! 화성에만 있는 특별한 신호를 만들어보면 어떨까? 지구와는 다른 화성만의 규칙!}
\sayalbi{삐릿♪ LED 3개가 동시에 깜빡이거나, 파도처럼 흐르는 신호를 만들 수 있어요!}
\end{storybox}

\subsubsection*{학습 목표}
\begin{itemize}[leftmargin=1.4em]
  \item 3개 LED 동시 빠른 깜빡임 "긴급 통과 신호"
  \item 파도처럼 순서대로 흐르는 "흐름 신호"
  \item 나만의 화성 신호 체계 완성
\end{itemize}

\begin{labbox}{샘플 코드 --- 화성 특별 신호 (파도 신호)}
\begin{minted}{python}
# 화성 특별 신호: 파도 신호
while True:
    for i in range(1, 4):       # 1→2→3
        led_on(i, 1.0)
        time.sleep(0.3)
        led_off(i)
    for i in range(3, 0, -1):  # 3→2→1
        led_on(i, 1.0)
        time.sleep(0.3)
        led_off(i)
\end{minted}
\end{labbox}

\begin{notebox}{확장 학습}
파도 신호는 \texttt{range(1,4)}와 \texttt{range(3,0,-1)}을 통한
\emph{인덱스 조작}을 처음 도입한다. LED 핀 번호를 변수 \texttt{i}로
받아 동적으로 제어한다는 발상은, 이후 차시에서 다수 LED를 다룰 때
핵심 패턴이 된다.
\end{notebox}
